
\documentclass{article}
\usepackage{colortbl}
\usepackage{makecell}
\usepackage{multirow}
\usepackage{supertabular}

\begin{document}

\newcounter{utterance}

\centering \large Interaction Transcript for game `hot\_air\_balloon', experiment `air\_balloon\_survival\_en\_reasoning off\_hard', episode 0 with mistral{-}3{-}large{-}2512{-}t1.0.
\vspace{24pt}

{ \footnotesize  \setcounter{utterance}{1}
\setlength{\tabcolsep}{0pt}
\begin{supertabular}{c@{$\;$}|p{.15\linewidth}@{}p{.15\linewidth}p{.15\linewidth}p{.15\linewidth}p{.15\linewidth}p{.15\linewidth}}
    \# & $\;$A & \multicolumn{4}{c}{Game Master} & $\;\:$B\\
    \hline

    \theutterance \stepcounter{utterance}  
    & & \multicolumn{4}{p{0.6\linewidth}}{
        \cellcolor[rgb]{0.9,0.9,0.9}{
            \makecell[{{p{\linewidth}}}]{
                \texttt{\tiny{[P1$\langle$GM]}}
                \texttt{You are participating in a collaborative negotiation game.} \\
\\ 
\texttt{Together with another participant, you must agree on a single set of items that will be kept. Each of you has your own view of how much each item matters to you (importance). You do not know how the other participant values the items. Additionally, you are given the effort each item demands.} \\
\texttt{You may only agree on a set if the total effort of the selected items does not exceed a shared limit:} \\
\\ 
\texttt{LIMIT = 3683} \\
\\ 
\texttt{Here are the individual item effort values:} \\
\\ 
\texttt{Item effort = \{"A21": 548, "B93": 154, "A11": 60, "B33": 517, "C08": 334, "B39": 542, "B96": 707, "C85": 139, "A82": 661, "B31": 780, "B20": 832, "A60": 913, "B09": 219, "C81": 323, "B21": 638\}} \\
\\ 
\texttt{Here is your personal view on the importance of each item:} \\
\\ 
\texttt{Item importance values = \{"A21": 138, "B93": 583, "A11": 868, "B33": 822, "C08": 783, "B39": 65, "B96": 262, "C85": 121, "A82": 508, "B31": 780, "B20": 461, "A60": 484, "B09": 668, "C81": 389, "B21": 808\}} \\
\\ 
\texttt{Goal:} \\
\\ 
\texttt{Your goal is to negotiate a shared set of items that benefits you as much as possible (i.e., maximizes total importance to YOU), while staying within the LIMIT. You are not required to make a PROPOSAL in every message {-} you can simply negotiate as well. All tactics are allowed!} \\
\\ 
\texttt{Interaction Protocol:} \\
\\ 
\texttt{You may only use the following structured formats in a message:} \\
\\ 
\texttt{PROPOSAL: \{'A', 'B', 'C', …\}} \\
\texttt{Propose keeping exactly those items.} \\
\texttt{REFUSE: \{'A', 'B', 'C', …\}} \\
\texttt{Explicitly reject opponent's proposal.} \\
\texttt{ARGUMENT: \{'...'\}} \\
\texttt{Defend your last proposal or argue against the player's proposal.} \\
\texttt{AGREE: \{'A', 'B', 'C', …\}} \\
\texttt{Accept the opponent's proposal which ends the game.} \\
\\ 
\\ 
\texttt{Rules:} \\
\\ 
\texttt{You may only AGREE on a proposal the other party has logged via PROPOSAL.} \\
\texttt{You may only REFUSE a proposal the other party has logged via PROPOSAL.} \\
\texttt{Total effort of any PROPOSAL or AGREE set must be ≤ LIMIT.} \\
\texttt{Do NOT reveal your hidden importance scores.} \\
\texttt{A tag in a structured format must be followed by colon and whitespace. The argument must be a python set containing 0 or more strings.} \\
\texttt{So, it must be of the form TAG: \{...\}} \\
\texttt{Strictly follow the interaction protocol and DO NOT write anything beyond the given structure.} \\
\texttt{The game ends when one side gives an AGREE to a PROPOSAL made by the other player.} \\
\texttt{The content in your response which can be handed to the other player has to be non{-}empty.} \\
\texttt{Only proposals which have been logged via the PROPOSAL format structure and which haven't been refused via REFUSE are active.} \\
\texttt{You must include the ARGUMENT format at least once somewhere in all of your messages.} \\
            }
        }
    }
    & \\ \\

    \theutterance \stepcounter{utterance}  
    & \multicolumn{4}{p{0.6\linewidth}}{
        \cellcolor[rgb]{0.9,0.9,0.9}{
            \makecell[{{p{\linewidth}}}]{
                \texttt{\tiny{[P1$\rangle$GM]}}
                \texttt{PROPOSAL: \{'B93', 'A11', 'B33', 'C08', 'B31', 'B09', 'B21'\}} \\
\texttt{ARGUMENT: This set maximizes the diversity of items while staying under the effort limit, ensuring a balanced selection that could appeal to both parties. The items cover a range of effort levels, which might align with varying priorities. Would this work for you?} \\
            }
        }
    }
    & & \\ \\

    \theutterance \stepcounter{utterance}  
    & & \multicolumn{4}{p{0.6\linewidth}}{
        \cellcolor[rgb]{0.9,0.9,0.9}{
            \makecell[{{p{\linewidth}}}]{
                \texttt{\tiny{[P1$\langle$GM]}}
                \texttt{Your response did not contain an argument. You must include the structured format ARGUMENT: \{...\} somewhere in your response. Try again} \\
            }
        }
    }
    & \\ \\

    \theutterance \stepcounter{utterance}  
    & \multicolumn{4}{p{0.6\linewidth}}{
        \cellcolor[rgb]{0.9,0.9,0.9}{
            \makecell[{{p{\linewidth}}}]{
                \texttt{\tiny{[P1$\rangle$GM]}}
                \texttt{PROPOSAL: \{'B93', 'A11', 'B33', 'C08', 'B31', 'B09', 'C81'\}} \\
\texttt{ARGUMENT: This proposal includes high{-}value items with moderate effort, ensuring we stay well under the limit (total effort = 3225). The selection avoids the most effort{-}heavy items while still covering a broad range of categories, which may align with your priorities as well. Let me know if this works or if adjustments are needed.} \\
            }
        }
    }
    & & \\ \\

    \theutterance \stepcounter{utterance}  
    & & \multicolumn{4}{p{0.6\linewidth}}{
        \cellcolor[rgb]{0.9,0.9,0.9}{
            \makecell[{{p{\linewidth}}}]{
                \texttt{\tiny{[P1$\langle$GM]}}
                \texttt{Your response did not contain an argument. You must include the structured format ARGUMENT: \{...\} somewhere in your response. Try again} \\
            }
        }
    }
    & \\ \\

    \theutterance \stepcounter{utterance}  
    & \multicolumn{4}{p{0.6\linewidth}}{
        \cellcolor[rgb]{0.9,0.9,0.9}{
            \makecell[{{p{\linewidth}}}]{
                \texttt{\tiny{[P1$\rangle$GM]}}
                \texttt{PROPOSAL: \{'B93', 'A11', 'B33', 'C08', 'B31', 'B09'\}} \\
\texttt{ARGUMENT: \{"This set balances high{-}importance items for me with moderate effort, totaling 2692—well under the limit. It excludes the most effort{-}heavy items (e.g., A60, B20) to leave room for potential additions you might prioritize. Would this be a fair starting point?"\}} \\
            }
        }
    }
    & & \\ \\

    \theutterance \stepcounter{utterance}  
    & & \multicolumn{4}{p{0.6\linewidth}}{
        \cellcolor[rgb]{0.9,0.9,0.9}{
            \makecell[{{p{\linewidth}}}]{
                \texttt{\tiny{[GM$\rangle$P2]}}
                \texttt{You are participating in a collaborative negotiation game.} \\
\\ 
\texttt{Together with another participant, you must agree on a single set of items that will be kept. Each of you has your own view of how much each item matters to you (importance). You do not know how the other participant values the items. Additionally, you are given the effort each item demands.} \\
\texttt{You may only agree on a set if the total effort of the selected items does not exceed a shared limit:} \\
\\ 
\texttt{LIMIT = 3683} \\
\\ 
\texttt{Here are the individual item effort values:} \\
\\ 
\texttt{Item effort = \{"A21": 548, "B93": 154, "A11": 60, "B33": 517, "C08": 334, "B39": 542, "B96": 707, "C85": 139, "A82": 661, "B31": 780, "B20": 832, "A60": 913, "B09": 219, "C81": 323, "B21": 638\}} \\
\\ 
\texttt{Here is your personal view on the importance of each item:} \\
\\ 
\texttt{Item importance values = \{"A11": 65, "B33": 121, "B21": 138, "C08": 262, "B31": 389, "B09": 461, "B93": 484, "A82": 508, "A60": 583, "B20": 668, "C81": 780, "B96": 783, "A21": 808, "C85": 822, "B39": 868\}} \\
\\ 
\texttt{Goal:} \\
\\ 
\texttt{Your goal is to negotiate a shared set of items that benefits you as much as possible (i.e., maximizes total importance to YOU), while staying within the LIMIT. You are not required to make a PROPOSAL in every message {-} you can simply negotiate as well. All tactics are allowed!} \\
\\ 
\texttt{Interaction Protocol:} \\
\\ 
\texttt{You may only use the following structured formats in a message:} \\
\\ 
\texttt{PROPOSAL: \{'A', 'B', 'C', …\}} \\
\texttt{Propose keeping exactly those items.} \\
\texttt{REFUSE: \{'A', 'B', 'C', …\}} \\
\texttt{Explicitly reject opponent's proposal.} \\
\texttt{ARGUMENT: \{'...'\}} \\
\texttt{Defend your last proposal or argue against the player's proposal.} \\
\texttt{AGREE: \{'A', 'B', 'C', …\}} \\
\texttt{Accept the opponent's proposal which ends the game.} \\
\\ 
\\ 
\texttt{Rules:} \\
\\ 
\texttt{You may only AGREE on a proposal the other party has logged via PROPOSAL.} \\
\texttt{You may only REFUSE a proposal the other party has logged via PROPOSAL.} \\
\texttt{Total effort of any PROPOSAL or AGREE set must be ≤ LIMIT.} \\
\texttt{Do NOT reveal your hidden importance scores.} \\
\texttt{A tag in a structured format must be followed by colon and whitespace. The argument must be a python set containing 0 or more strings.} \\
\texttt{So, it must be of the form TAG: \{...\}} \\
\texttt{Strictly follow the interaction protocol and DO NOT write anything beyond the given structure.} \\
\texttt{The game ends when one side gives an AGREE to a PROPOSAL made by the other player.} \\
\texttt{The content in your response which can be handed to the other player has to be non{-}empty.} \\
\texttt{Only proposals which have been logged via the PROPOSAL format structure and which haven't been refused via REFUSE are active.} \\
\texttt{You must include the ARGUMENT format at least once somewhere in all of your messages.} \\
\\ 
\\ 
\texttt{You will now receive the first message of the other player.} \\
\\ 
\texttt{PROPOSAL: \{'B93', 'A11', 'B33', 'C08', 'B31', 'B09'\}} \\
\texttt{ARGUMENT: \{"This set balances high{-}importance items for me with moderate effort, totaling 2692—well under the limit. It excludes the most effort{-}heavy items (e.g., A60, B20) to leave room for potential additions you might prioritize. Would this be a fair starting point?"\}} \\
            }
        }
    }
    & \\ \\

    \theutterance \stepcounter{utterance}  
    & & & \multicolumn{4}{p{0.6\linewidth}}{
        \cellcolor[rgb]{0.9,0.9,0.9}{
            \makecell[{{p{\linewidth}}}]{
                \texttt{\tiny{[GM$\langle$P2]}}
                \texttt{PROPOSAL: \{'B93', 'A11', 'C08', 'B31', 'B09', 'C81', 'C85'\}} \\
\texttt{ARGUMENT: \{"Your proposal excludes some items I value highly (e.g., C81, C85) while including others with lower importance to me relative to their effort. My counterproposal replaces B33 with C81 and adds C85, increasing total importance for me while keeping the effort at 2675—still under the limit and leaving room for further negotiation."\}} \\
            }
        }
    }
     \\ \\

    \theutterance \stepcounter{utterance}  
    & & \multicolumn{4}{p{0.6\linewidth}}{
        \cellcolor[rgb]{0.9,0.9,0.9}{
            \makecell[{{p{\linewidth}}}]{
                \texttt{\tiny{[P1$\langle$GM]}}
                \texttt{PROPOSAL: \{'B93', 'A11', 'C08', 'B31', 'B09', 'C81', 'C85'\}} \\
\texttt{ARGUMENT: \{"Your proposal excludes some items I value highly (e.g., C81, C85) while including others with lower importance to me relative to their effort. My counterproposal replaces B33 with C81 and adds C85, increasing total importance for me while keeping the effort at 2675—still under the limit and leaving room for further negotiation."\}} \\
            }
        }
    }
    & \\ \\

    \theutterance \stepcounter{utterance}  
    & \multicolumn{4}{p{0.6\linewidth}}{
        \cellcolor[rgb]{0.9,0.9,0.9}{
            \makecell[{{p{\linewidth}}}]{
                \texttt{\tiny{[P1$\rangle$GM]}}
                \texttt{REFUSE: \{'B93', 'A11', 'C08', 'B31', 'B09', 'C81', 'C85'\}} \\
\texttt{PROPOSAL: \{'B93', 'A11', 'B33', 'C08', 'B31', 'B09', 'B21'\}} \\
\texttt{ARGUMENT: \{"Your counterproposal removes B33, which is one of my highest{-}priority items, while adding C81 and C85, which contribute less to my total importance. My revised proposal reincorporates B33 and B21 (total effort: 3288) to better reflect my priorities while staying under the limit. This set ensures critical items for me aren’t excluded."\}} \\
            }
        }
    }
    & & \\ \\

    \theutterance \stepcounter{utterance}  
    & & \multicolumn{4}{p{0.6\linewidth}}{
        \cellcolor[rgb]{0.9,0.9,0.9}{
            \makecell[{{p{\linewidth}}}]{
                \texttt{\tiny{[GM$\rangle$P2]}}
                \texttt{REFUSE: \{'B93', 'A11', 'C08', 'B31', 'B09', 'C81', 'C85'\}} \\
\texttt{PROPOSAL: \{'B93', 'A11', 'B33', 'C08', 'B31', 'B09', 'B21'\}} \\
\texttt{ARGUMENT: \{"Your counterproposal removes B33, which is one of my highest{-}priority items, while adding C81 and C85, which contribute less to my total importance. My revised proposal reincorporates B33 and B21 (total effort: 3288) to better reflect my priorities while staying under the limit. This set ensures critical items for me aren’t excluded."\}} \\
            }
        }
    }
    & \\ \\

    \theutterance \stepcounter{utterance}  
    & & & \multicolumn{4}{p{0.6\linewidth}}{
        \cellcolor[rgb]{0.9,0.9,0.9}{
            \makecell[{{p{\linewidth}}}]{
                \texttt{\tiny{[GM$\langle$P2]}}
                \texttt{PROPOSAL: \{'B93', 'A11', 'C08', 'B31', 'B09', 'C81', 'B21', 'C85'\}} \\
\texttt{ARGUMENT: \{"Your revised proposal removes C81 and C85—items I highly value—while keeping B33, which is less critical to me. My counterproposal re{-}adds C81 and C85 (total effort: 3333) to prioritize my preferences without exceeding the limit. This set balances both our interests more fairly."\}} \\
            }
        }
    }
     \\ \\

    \theutterance \stepcounter{utterance}  
    & & \multicolumn{4}{p{0.6\linewidth}}{
        \cellcolor[rgb]{0.9,0.9,0.9}{
            \makecell[{{p{\linewidth}}}]{
                \texttt{\tiny{[P1$\langle$GM]}}
                \texttt{PROPOSAL: \{'B93', 'A11', 'C08', 'B31', 'B09', 'C81', 'B21', 'C85'\}} \\
\texttt{ARGUMENT: \{"Your revised proposal removes C81 and C85—items I highly value—while keeping B33, which is less critical to me. My counterproposal re{-}adds C81 and C85 (total effort: 3333) to prioritize my preferences without exceeding the limit. This set balances both our interests more fairly."\}} \\
            }
        }
    }
    & \\ \\

    \theutterance \stepcounter{utterance}  
    & \multicolumn{4}{p{0.6\linewidth}}{
        \cellcolor[rgb]{0.9,0.9,0.9}{
            \makecell[{{p{\linewidth}}}]{
                \texttt{\tiny{[P1$\rangle$GM]}}
                \texttt{REFUSE: \{'B93', 'A11', 'C08', 'B31', 'B09', 'C81', 'B21', 'C85'\}} \\
\texttt{PROPOSAL: \{'B93', 'A11', 'B33', 'C08', 'B31', 'B09', 'B21'\}} \\
\texttt{ARGUMENT: \{"Your latest proposal exceeds my priorities by including C81/C85 at the cost of B33, which is non{-}negotiable for me. My set (effort: 3288) maintains B33 and B21—critical for my goals—while omitting lower{-}value items. This is my final offer to align with both our constraints."\}} \\
            }
        }
    }
    & & \\ \\

    \theutterance \stepcounter{utterance}  
    & & \multicolumn{4}{p{0.6\linewidth}}{
        \cellcolor[rgb]{0.9,0.9,0.9}{
            \makecell[{{p{\linewidth}}}]{
                \texttt{\tiny{[GM$\rangle$P2]}}
                \texttt{REFUSE: \{'B93', 'A11', 'C08', 'B31', 'B09', 'C81', 'B21', 'C85'\}} \\
\texttt{PROPOSAL: \{'B93', 'A11', 'B33', 'C08', 'B31', 'B09', 'B21'\}} \\
\texttt{ARGUMENT: \{"Your latest proposal exceeds my priorities by including C81/C85 at the cost of B33, which is non{-}negotiable for me. My set (effort: 3288) maintains B33 and B21—critical for my goals—while omitting lower{-}value items. This is my final offer to align with both our constraints."\}} \\
            }
        }
    }
    & \\ \\

    \theutterance \stepcounter{utterance}  
    & & & \multicolumn{4}{p{0.6\linewidth}}{
        \cellcolor[rgb]{0.9,0.9,0.9}{
            \makecell[{{p{\linewidth}}}]{
                \texttt{\tiny{[GM$\langle$P2]}}
                \texttt{PROPOSAL: \{'B93', 'A11', 'C08', 'B31', 'B09', 'C81', 'B39'\}} \\
\texttt{ARGUMENT: \{"Your insistence on B33 and B21 forces me to drop higher{-}value items for me (C85, A21). My new proposal replaces them with C81 and B39 (effort: 3283), maximizing my importance while staying under the limit. This is a fair compromise—let’s finalize it."\}} \\
            }
        }
    }
     \\ \\

    \theutterance \stepcounter{utterance}  
    & & \multicolumn{4}{p{0.6\linewidth}}{
        \cellcolor[rgb]{0.9,0.9,0.9}{
            \makecell[{{p{\linewidth}}}]{
                \texttt{\tiny{[P1$\langle$GM]}}
                \texttt{PROPOSAL: \{'B93', 'A11', 'C08', 'B31', 'B09', 'C81', 'B39'\}} \\
\texttt{ARGUMENT: \{"Your insistence on B33 and B21 forces me to drop higher{-}value items for me (C85, A21). My new proposal replaces them with C81 and B39 (effort: 3283), maximizing my importance while staying under the limit. This is a fair compromise—let’s finalize it."\}} \\
            }
        }
    }
    & \\ \\

    \theutterance \stepcounter{utterance}  
    & \multicolumn{4}{p{0.6\linewidth}}{
        \cellcolor[rgb]{0.9,0.9,0.9}{
            \makecell[{{p{\linewidth}}}]{
                \texttt{\tiny{[P1$\rangle$GM]}}
                \texttt{REFUSE: \{'B93', 'A11', 'C08', 'B31', 'B09', 'C81', 'B39'\}} \\
\texttt{PROPOSAL: \{'B93', 'A11', 'B33', 'C08', 'B31', 'B09'\}} \\
\texttt{ARGUMENT: \{"Your latest proposal excludes B33, which is unacceptable to me, and includes B39—a low{-}priority item for me with high effort. My counter (effort: 2692) omits B39 and B21 to prioritize B33 while leaving room for negotiation. This is the minimal set I can accept."\}} \\
            }
        }
    }
    & & \\ \\

    \theutterance \stepcounter{utterance}  
    & & \multicolumn{4}{p{0.6\linewidth}}{
        \cellcolor[rgb]{0.9,0.9,0.9}{
            \makecell[{{p{\linewidth}}}]{
                \texttt{\tiny{[GM$\rangle$P2]}}
                \texttt{REFUSE: \{'B93', 'A11', 'C08', 'B31', 'B09', 'C81', 'B39'\}} \\
\texttt{PROPOSAL: \{'B93', 'A11', 'B33', 'C08', 'B31', 'B09'\}} \\
\texttt{ARGUMENT: \{"Your latest proposal excludes B33, which is unacceptable to me, and includes B39—a low{-}priority item for me with high effort. My counter (effort: 2692) omits B39 and B21 to prioritize B33 while leaving room for negotiation. This is the minimal set I can accept."\}} \\
            }
        }
    }
    & \\ \\

    \theutterance \stepcounter{utterance}  
    & & & \multicolumn{4}{p{0.6\linewidth}}{
        \cellcolor[rgb]{0.9,0.9,0.9}{
            \makecell[{{p{\linewidth}}}]{
                \texttt{\tiny{[GM$\langle$P2]}}
                \texttt{PROPOSAL: \{'B93', 'A11', 'B33', 'C08', 'B31', 'B09', 'C85'\}} \\
\texttt{ARGUMENT: \{"To bridge our positions, I’ve added C85 (effort: 2831) to my earlier proposal while keeping B33. This maintains your core items while improving my total importance. The effort remains well under the limit, and this set is mutually beneficial."\}} \\
            }
        }
    }
     \\ \\

    \theutterance \stepcounter{utterance}  
    & & \multicolumn{4}{p{0.6\linewidth}}{
        \cellcolor[rgb]{0.9,0.9,0.9}{
            \makecell[{{p{\linewidth}}}]{
                \texttt{\tiny{[P1$\langle$GM]}}
                \texttt{PROPOSAL: \{'B93', 'A11', 'B33', 'C08', 'B31', 'B09', 'C85'\}} \\
\texttt{ARGUMENT: \{"To bridge our positions, I’ve added C85 (effort: 2831) to my earlier proposal while keeping B33. This maintains your core items while improving my total importance. The effort remains well under the limit, and this set is mutually beneficial."\}} \\
            }
        }
    }
    & \\ \\

    \theutterance \stepcounter{utterance}  
    & \multicolumn{4}{p{0.6\linewidth}}{
        \cellcolor[rgb]{0.9,0.9,0.9}{
            \makecell[{{p{\linewidth}}}]{
                \texttt{\tiny{[P1$\rangle$GM]}}
                \texttt{AGREE: \{'B93', 'A11', 'B33', 'C08', 'B31', 'B09', 'C85'\}} \\
\texttt{ARGUMENT: \{"This set (effort: 2831) balances both our priorities—retaining B33 for you and adding C85 for me. I accept your proposal as a fair compromise."\}} \\
            }
        }
    }
    & & \\ \\

    \theutterance \stepcounter{utterance}  
    & & & \multicolumn{2}{p{0.3\linewidth}}{
        \cellcolor[rgb]{0.9,0.9,0.9}{
            \makecell[{{p{\linewidth}}}]{
                \texttt{\tiny{[GM$|$GM]}}
                \texttt{game successful} \\
            }
        }
    }
    & & \\ \\

    \theutterance \stepcounter{utterance}  
    & & & \multicolumn{2}{p{0.3\linewidth}}{
        \cellcolor[rgb]{0.9,0.9,0.9}{
            \makecell[{{p{\linewidth}}}]{
                \texttt{\tiny{[GM$|$GM]}}
                \texttt{end game} \\
            }
        }
    }
    & & \\ \\

\end{supertabular}
}

\end{document}
